\par\noindent\textbf{(1)} 取左侧y1=e\textasciicircum{}\{i(k1x−ωt)\}+r e\textasciicircum{}\{i(−k1x−ωt)\}、右侧y2=t e\textasciicircum{}\{i(k2x−ωt)\}，k\_j=ω/c\_j。界面位移连续给1+r=t，张力T∂xy连续给k1(1−r)=k2t；这两个条件来自质量零的界面，不应另外添加惯性跳跃项。
\par\noindent\textbf{(2)} 联立两条件得r=(k1−k2)/(k1+k2)、t=2k1/(k1+k2)。定义弦特性机械阻抗Z\_j=T/c\_j=√(Tμ\_j)，因k1/k2=Z1/Z2，故r=(Z1−Z2)/(Z1+Z2)、t=2Z1/(Z1+Z2)。当Z2>Z1，位移反射符号为负。
\par\noindent\textbf{(3)} 行波时间平均功率P\_j=Z\_jω²|A\_j|²/2，所以反射功率比R=|r|²、透射比𝒯=(Z2/Z1)|t|²=4Z1Z2/(Z1+Z2)²。直接代数求和R+𝒯=[(Z1−Z2)²+4Z1Z2]/(Z1+Z2)²=1，满足无耗散界面能量守恒。
\par\noindent\textbf{(4)} 当Z2≫Z1时r→−1、t→0，呈近固定端；当Z2≪Z1时r→1、t→2，但功率透射因Z2/Z1趋零而消失。若Z1=Z2，r=0且𝒯=1，即使两段弦的密度或几何不同，只要行波阻抗相等就不会发生反射。
